\documentclass[11pt,a4paper]{article}
\usepackage[utf8]{inputenc}
\usepackage[spanish]{babel}
\usepackage{amsmath}
\usepackage{amsfonts}
\usepackage{amssymb}
\usepackage{booktabs}
\usepackage{array}
\usepackage{geometry}
\geometry{top=2.5cm, bottom=2.5cm, left=2cm, right=2cm}

\title{Ecuaciones Operativas - Subsistema de Entrada (M/M/2)}
\author{Cátedra de Investigación de Operaciones / Teoría de Colas}
\date{\today}

\begin{document}

\maketitle

A continuación se presentan las ecuaciones operativas correspondientes al **Subsistema Entrada ($M/M/2$)**, estructuradas para su uso directo en reportes y análisis académicos.

\vspace{0.5cm}

\begin{table}[h!]
\centering
\small
\renewcommand{\arraystretch}{2.2} % Espaciado vertical para las fracciones
\begin{tabular}{c >{\displaystyle}l p{8cm}}
\toprule
\textbf{\#} & \multicolumn{1}{c}{\textbf{Ecuación}} & \textbf{Descripción del Parámetro (Cátedra)} \\
\midrule
1.1 & \rho = \frac{\lambda}{2\mu} & Tasa de ocupación (Factor de utilización del sistema) \\
\midrule
1.2 & P_0 = \left[ 1 + 2\rho + \frac{(2\rho)^2}{2!(1 - \rho)} \right]^{-1} & Probabilidad de que el sistema esté vacío \\
\midrule
1.3 & P_n = \frac{(2\rho)^n}{n!} \cdot P_0 \quad (\text{para } n < 2) & Probabilidad de que $n$ entidades estén en el sistema \\
\midrule
1.4 & P_n = 2 \cdot \rho^n \cdot P_0 \quad (\text{para } n \ge 2) & Probabilidad de que $n$ entidades estén en el sistema \\
\midrule
1.5 & P_K = 0 & Probabilidad de bloqueo (Capacidad infinita) \\
\midrule
1.6 & \lambda_{ef} = \lambda & Tasa efectiva de llegada al subsistema \\
\midrule
1.7 & L_q = \left[ \frac{2\rho^2}{(1 - \rho)^2} \right] \cdot P_0 & Número de entidades en la cola del sistema \\
\midrule
1.8 & L_w = L_q + \frac{\lambda}{\mu} & Número esperado de entidades en el sistema de colas \\
\midrule
1.9 & W_q = \frac{L_q}{\lambda} & Tiempo promedio que una entidad espera en la cola \\
\midrule
1.10 & W_w = W_q + \frac{1}{\mu} & Tiempo promedio que una entidad espera en el sistema \\
\midrule
1.11 & P(W_q > t) = \left[ \frac{(2\rho)^2}{2!(1 - \rho)} \right] \cdot P_0 \cdot e^{-2\mu(1 - \rho)t} & Probabilidad de que $W_q$ sea mayor que $t$ \\
\midrule
1.12 & P(W_w > t) = e^{-\mu t} \cdot \left[ 1 + \left( \frac{(2\rho)^2}{2!(1 - \rho)} \right) \cdot P_0 \right] & Probabilidad de que $W_w$ sea mayor que $t$ \\
\bottomrule
\end{tabular}
\caption{Ecuaciones Operativas para el modelo $M/M/2$.}
\label{tab:ecuaciones_operativas}
\end{table}

\end{document}
